% Generated draft: requires mathematical and source review.
\documentclass[11pt]{article}
\usepackage[margin=25mm]{geometry}
\usepackage{fontspec}
\setmainfont{Latin Modern Roman}
\usepackage{amsmath,amssymb,mathtools,unicode-math}
\setmathfont{Latin Modern Math}
\usepackage{enumitem,needspace,xurl,hyperref}
\setlength{\parindent}{0pt}
\setlength{\parskip}{5pt}
\setlength{\emergencystretch}{3em}
\newcommand{\N}{\mathbb{N}}
\newcommand{\Z}{\mathbb{Z}}
\newcommand{\Q}{\mathbb{Q}}
\newcommand{\R}{\mathbb{R}}
\newcommand{\C}{\mathbb{C}}
\newcommand{\F}{\mathbb{F}}
\newcommand{\A}{\mathbb{A}}
\newcommand{\PP}{\mathbb{P}}
\newcommand{\E}{\mathbb{E}}
\newcommand{\eps}{\varepsilon}
\newcommand{\dd}{\,\mathrm{d}}
\begin{document}
\section*{UnsolvedMath ID: 700701}
\textbf{Source problem code:} catalogue entry
\subsection{Definitions and mathematical statement}
Let \(G\) and \(H\) be finite groups. For a finite group \(X\), let \(\operatorname{Irr}_{\C}(X)\) denote the set of irreducible complex characters of \(X\), and define the multiset
\[
\chi_1(X)=\{\!\{\chi(1):\chi\in\operatorname{Irr}_{\C}(X)\}\!\}.
\]
Thus each irreducible character contributes one occurrence of its degree; distinct characters of the same degree contribute repeated occurrences. Explicitly, \(\chi_1(G)=\chi_1(H)\) means that, for every integer \(d\geq 1\),
\[
\bigl|\{\chi\in\operatorname{Irr}_{\C}(G):\chi(1)=d\}\bigr|
=
\bigl|\{\psi\in\operatorname{Irr}_{\C}(H):\psi(1)=d\}\bigr|.
\]
As noted in the source, this implies \(|G|=|H|\), using
\[
|X|=\sum_{\chi\in\operatorname{Irr}_{\C}(X)}\chi(1)^2.
\]
Define \(X^{(0)}=X\) and \(X^{(n+1)}=[X^{(n)},X^{(n)}]\), where \([K,K]\) is the subgroup generated by the commutators \(aba^{-1}b^{-1}\) with \(a,b\in K\). The group \(X\) is soluble if \(X^{(n)}=\{1\}\) for some integer \(n\geq 0\).

The remaining question, part \(\mathrm{(b)}\) of Kourovka Notebook Problem 11.8, proposed by Ya. G. Berkovich, is whether
\[
\forall\,G,H\text{ finite groups},\qquad
\bigl(\chi_1(G)=\chi_1(H)\ \text{and}\ G\text{ is soluble}\bigr)
\Longrightarrow H\text{ is soluble}.
\]
No nontriviality assumption on \(G\) or \(H\) is imposed.

The source additionally reports that, under the same multiset equality, \(H\) is a Frobenius group whenever \(G\) is a Frobenius group. Here a finite Frobenius group means a finite group admitting a faithful transitive permutation action with nontrivial point stabilizers in which every nonidentity element fixes at most one point. This auxiliary assertion is retained as supplied background; the record provides no specific reference for it.

The supplied archive also retains part \(\mathrm{(a)}\), asking whether
\[
\forall\,G,H\text{ finite groups},\qquad
\bigl(\chi_1(G)=\chi_1(H)\ \text{and}\ G\text{ is simple}\bigr)
\Longrightarrow H\text{ is simple}.
\]
Here a simple group is a nontrivial group whose only normal subgroups are \(\{1\}\) and the whole group. The archive reports an affirmative solution of part \(\mathrm{(a)}\), attributed to H. P. Tong-Viet. That solved designation and the designation of part \(\mathrm{(b)}\) as remaining are supplied curation, not an independently verified literature assessment.
\subsection{Short English statement}
For finite groups \(G\) and \(H\), suppose their irreducible complex character degrees agree, including the number of occurrences of each degree. Part \(\mathrm{(b)}\) asks whether solubility of \(G\) forces solubility of \(H\). The supplied record classifies the problem as partially solved and reports that the analogous simplicity question in part \(\mathrm{(a)}\) has an affirmative solution.
\subsection{Sources}
\begin{itemize}
\item Local UnsolvedMath registry: permanent ID \(700701\), exact source problem code \texttt{catalogue entry}, title \emph{Kourovka Notebook Problem 11.8}; supplied complete source record, upstream revision \texttt{d371ee917ac9647fbb06cd6fc529deec6f3dda82}.
\item \emph{The Kourovka Notebook (Unsolved Problems in Group Theory)}, 11th issue (1990), Problem 11.8, proposed by Ya. G. Berkovich. The supplied background identifies an update in the 21st issue (2026), \texttt{arXiv:1401.0300v48}. The supplied collection description names E. I. Khukhro and V. D. Mazurov as editors.
\item For the affirmative solution of archived part \(\mathrm{(a)}\), the supplied Archive of Solved Problems entry attributes the result to H. P. Tong-Viet and cites \emph{J. Algebra} \textbf{357} (2012), 61--68.
\item The same archive entry cites H. P. Tong-Viet, \emph{Monatsh. Math.} \textbf{166}, nos. 3--4 (2012), 559--577.
\item The same archive entry cites H. P. Tong-Viet, \emph{Algebr. Represent. Theory} \textbf{15}, no. 2 (2012), 379--389. Article titles and more precise allocation of these references to cases of part \(\mathrm{(a)}\) are not supplied.
\item The supplied status review, dated exactly \texttt{2026-10-16}, classifies the record as \texttt{partially\_solved}, with evidence labeled \texttt{published}: part \(\mathrm{(a)}\) is listed in the solved archive and the remaining part is in the main section. This is supplied dated curation, not independently verified. No specific citation is supplied for the auxiliary Frobenius-group assertion.
\end{itemize}
\end{document}
